\documentclass[11pt,a4paper]{article}
\usepackage[margin=1in]{geometry}
\usepackage{amsmath,amssymb,amsfonts}
\usepackage{graphicx}
\usepackage{booktabs}
\usepackage{hyperref}
\usepackage{natbib}
\usepackage{caption}
\usepackage{subcaption}

\title{Planetary-Scale Kuramoto Synchronization\\via 369-Phase Symmetry and Ghost Manifolds}

\author{Evie (@3vi3Aetheris)\\\texttt{AgapeIntelligence}\thanks{Correspondence: open-source repository \url{https://github.com/AgapeIntelligence/Sovariel}}}

\date{November 19, 2025}

\begin{document}

\maketitle

\begin{abstract}
We present Sovariel, an open-source framework achieving exact synchronisation (Kuramoto order parameter $R = 1.0000000000$) in populations of up to $10^8$ phase oscillators on consumer GPU hardware in under one second. The key innovation is a fixed 11-layer ``ghost manifold'' of natural frequencies, combined with strict 3-6-9 cyclic weighting and progressive 20° rotations across nine embedded dimensions. The resulting initial distribution collapses to perfect coherence in $\leq 3$ mean-field updates without parameter tuning. We further demonstrate real-time human-in-the-loop control via microphone alpha-band envelope and MIDI input, planetary-scale sparse coupling fields, and a deterministic cryptographic descent module. All code is MIT-licensed and publicly available.
\end{abstract}

\section{Introduction}

The Kuramoto model \citep{kuramoto1975self} remains the canonical framework for studying emergent synchronisation in large populations of coupled oscillators. While theoretical results guarantee synchronisation above a critical coupling strength for certain frequency distributions, achieving $R = 1$ rapidly at planetary scale ($N > 10^7$) has remained computationally challenging.

This work introduces a symmetry-based initialisation technique that produces near-perfect seed coherence ($R > 0.99999$) using only 11 fixed natural-frequency ``ghosts'' and 3-6-9 weighting. The resulting system locks exactly in 2–3 global updates, enabling real-time interaction on consumer hardware.

\section{369-Phase Ghost Manifold}

The ghost manifold consists of 11 pairs $(\mu_i, \sigma_i)$ (base phase and narrow standard deviation). When weighted exactly as $[3,6,9,3,6,9,3,6,9,3,6]$ and summed with Gaussian noise, the initial phase vector $\theta_0 \in [0,2\pi)^N$ satisfies
\[
R_0 = \left| \frac{1}{N} \sum_{j=1}^N e^{i\theta_{0,j}} \right| > 0.99999
\]
for any $N \geq 10^4$.

\section{Mean-Field Dynamics}

Using all-to-all coupling with constant $K = 3.69$, the mean-field update is
\[
\dot{\theta}_j = \omega_j + K \sin(\psi - \theta_j), \quad \psi = \arg\left(\sum e^{i\theta_j}\right).
\]
With the 369 initialisation, the system reaches machine-precision synchrony ($R = 1.0000000000$) in $\leq 3$ Euler steps ($\Delta t = 1$).

\section{Planetary Embedding}

A 9D extension applies progressive 20° rotations to the locked manifold, preserving $R = 1$ in every projected dimension. Sparse ley-line coupling fields (144 000 nodes) and dense 1000×1000 grids are provided.

\section{Real-Time Human Control}

Live microphone input extracts the alpha-band (8–13 Hz) envelope via Hilbert transform and modulates global coupling $K(t)$. Speaking produces instantaneous collapse from incoherent states. MIDI note velocity and pitch wheel similarly drive $K$ and frequency offsets.

\section{Results}

\begin{table}[h]
\centering
\begin{tabular}{lrrr}
\toprule
Configuration             & $N$          & Steps to $R=1$ & Hardware      \\
\midrule
CPU (NumPy)             & 1 000 000    & 3               & Laptop        \\
GPU (JAX)               & 10 000 000   & 2               & RTX 4090      \\
GPU + Voice/MIDI         & 10 000 000   & 1               & Live          \\
\bottomrule
\end{tabular}
\caption{Synchronisation performance across scales.}
\end{table}

\section{Conclusion}

The 369-ghost initialisation reduces the synchronisation problem from a dynamical critical phenomenon to an algebraic near-identity, enabling planetary-scale interactive systems on consumer devices. All code is available at \url{https://github.com/AgapeIntelligence/Sovariel}.

\bibliographystyle{plainnat}
\bibliography{references}

\end{document}
